\documentclass[UTF8,fontset=none,11pt,a4paper]{ctexart}
\usepackage{ipara-handbook}
\handbooksetup{DS-12 外部排序}{408 · 数据结构 DS-12}{Block I/O · Runs · Multiway Merge · Passes}
\begin{document}
\handbooktitle{外部排序}{数据不能全进内存时，怎样减少块 I/O 轮数}{408 · 数据结构 Topic DS-12}

\begin{corebox}{母问题：当 I/O 远贵于比较，排序模型怎样改写？}
外部排序的生成链是
\[
\text{Block-I/O Cost}\to\text{Initial Runs}\to\text{Multiway Merge}\to\text{Pass Count}\to\text{I/O Tradeoff}.
\]
数据先分成能装入内存的初始归并段并在内存中排序，再用 $k$ 路归并逐轮合并。总成本主要由读写块次数和归并轮数决定，不能直接套内部排序的比较次数。
\end{corebox}
\begin{methodbox}{本册定位与 Owner 边界}
\begin{itemize}
\item \kw{Owns：}初始归并段、置换选择、多路归并、败者树、缓冲区、I/O 轮数与终止条件。
\item \kw{Uses：}Atlas Foundation 的 cost vector；存储系统的 block 概念只作为外部接口，不复制文件系统机制。
\item \kw{Stops：}内存内局部排序归 DS11；设备调度、缓存一致性和文件系统完整语义不在本册。
\end{itemize}
\end{methodbox}
\tableofcontents\newpage

\section{初始段与多路归并}
设内存可容纳 $M$ 条记录，先读入至多 $M$ 条并排序，写成一个 run。归并时每个输入 run 保留一个当前首元素，取最小者输出并补充同一路的下一条；候选结构可用最小堆或败者树。若有 $r$ 个初始段、每轮 $k$ 路，轮数约为 $\lceil\log_k r\rceil$，每轮都要读写全部数据。
\begin{examplebox}{为什么败者树值得出现}
多路归并每输出一条记录都要更新“哪一路最小”。普通堆和败者树都维护候选极值；败者树显式保存比较路径，便于重复更新同一路。二者都不改变多路归并的正确性不变量。
\end{examplebox}

\section{I/O 不变量与终止}
每个 run 的当前指针只向前移动；输出序列始终是所有当前首元素中的最小值；某一路耗尽后从候选集合移除；所有输入耗尽才结束。缓冲区满/空是 I/O 事件，不是记录顺序事件，必须分开书写。块大小、内存缓冲数量和归并路数改变的是轮数与常数，不改变“每一路有序”的前提。

\section{代码与测试证据}
完整实现位于 \codepath{code/ds12_external_sort.hpp}，测试位于 \codepath{code/ds12_external_sort_test.cpp}。可执行核心用内存中的 run 模拟外部边界，输出合并结果并计算所需轮数。
\lstinputlisting[language=C++,firstline=10,lastline=82,firstnumber=10,numbers=left,caption={有序 runs 的多路归并与轮数估算}]{30_408/10_数据结构/12_外部排序/code/ds12_external_sort.hpp}
\lstinputlisting[language=C++,firstline=8,lastline=40,firstnumber=8,numbers=left,caption={空 run、单 run 与多路耗尽测试}]{30_408/10_数据结构/12_外部排序/code/ds12_external_sort_test.cpp}
\begin{lstlisting}[language=bash]
clang++ -std=c++17 -Wall -Wextra -Werror -pedantic code/ds12_external_sort_test.cpp -o /tmp/ds12_test
/tmp/ds12_test
\end{lstlisting}

\section{复杂度与概念边界}
\begin{center}\small
\begin{tabularx}{\linewidth}{L{2.8cm}C{1.8cm}C{1.8cm}YY}\toprule
阶段 & 比较/更新 & I/O 视角 & 不变量 & 边界\\\midrule
初始段 & 每段内存排序 & 读写 $N$ & 每段有序 & $N\le M$\\
$k$ 路归并 & 每条更新候选 & 每轮读写 $N$ & 各路首元素候选 & 空 run\\
总过程 & 依赖实现 & $O(N\log_k r)$ 轮次量级 & 指针单调前进 & 输入耗尽\\\bottomrule\end{tabularx}\end{center}

\section{做题控制协议}
先标 $N,M,k,r$，再算初始段数量和归并轮数；每一轮分别写“读、比较、输出、补位、写”。看到全部数据能进内存，停止套外部 I/O 公式并转 DS11；看到块大小或缓冲区变化，重新计算轮数而非只比较关键字次数。
\begin{corebox}{压缩复原}
五问：内存一次容纳多少？初始 run 怎样形成？当前候选来自几路？哪一路耗尽如何退出？总成本按比较还是块 I/O 计？
\end{corebox}
\end{document}
